%! root=../main.tex
\section{Discussion and Future Work}~\label{sec:disc}

\ixb{Embarrassing first draft is fine!}

mainly about (apparent) limitations

\begin{itemize}
    \item Lifter is handwritten and trusted; mitigated by related work (verified lifting)
    \item Extraction to OCaml, trusted, uses ints; mitigated by related work (verified extraction) and future work
    \item ELF Packager is trusted; open problem (dangerous!), requires a formally verified loader
    \item does not protect from srop
    \item protection limited to binaries that make no system calls, rare
    \item assumes no concurrency or shared memory, rare
\end{itemize}

\pname takes a big step towards high assurance binary defense by formally verifying the binary rewriter, but it still trusts other components of the system.
Underlying the rewriter is the internal lifter, a manual implementation of the ARMv7 ISA semantics as a Gallina function translating code bytes to PIL statements.
Creating trustworthy mechanized semantics of ISAs is a subject of ongoing research~\cite{arm}

While \pname represents a big step towards high assurance binary defenses, there are still many weaknesses left to address.
First, the entire work rests on our trusted internal lifter, the Gallina function translating code bytes to PIL statements.
Verified and trustworthy lifting are subjects of ongoing research, but verifying PIL lifters is an open problem.
Second, the extraction is trusted.
The OCaml extractor is verified~\wc{trustworthy}, but in our extraction we extracted Rocq's unbounded integer type \texttt{Z} to OCaml's bounded \texttt{int}.
This opens the door to integer overflow and underflow bugs in the extracted code, but it makes the rewriter usable.
Extracting it without this optimization, that is an OCaml variant \texttt{Z} with the same structure as Rocq's, results in XXx longer rewrite times.
Rocq offers the capability of proving that this optimization is sound, i.e., that no over/underflows are possible, but we leave it to future work/industrial adoption.
Third, the ELF Packager is beyond the scope of formal verification because any specification of its correctness relies on the behavior of the loader.
To bring that part of the system into scope requires the availalability of a formally verified loader.


\todo{Future work
\begin{itemize}
    \item call-return semantics for trace properties
    \item call-return semantics for overlapping code?
    \item total correctness?
    \item position independent proofs?
\end{itemize}
}

\subsection{Generalizability}
- beyond armv7 to armv8 (maybe)
- beyond RISC to CISC

(scope aspect) formal verification of
    - decoder
    - merger (???)

\subsection{Coverage}
\comment{%
- **Limited scope of verification.** Current work focuses only on the `.text` section (ARMv7 instruction streams), leaving other components of the binary lifecycle unverified.

- **Unverified dependencies in end-to-end rewriting.** Extending coverage to the full pipeline, such as ELF packaging and program loading, requires external tools whose correctness we do not verify.
  - For example, the rewriting pipeline relies on libraries such as `lief` for ELF manipulation.
  - These dependencies fall outside the formal model, weakening end-to-end guarantees.

- **Challenge of full-stack formalization.** Rewriting and verifying the entire dependency chain in a formally verifiable language would be prohibitively complex.
  - Formal verification of systems implemented in general-purpose languages remains difficult and resource-intensive.

- **Alternative design direction.** One could instead formally verify the program loader against a well-defined security specification.
  - However, this approach also introduces significant complexity and remains a challenging open problem.
}